% !TEX root = ../main.tex
%% Corrigés — Chapitre 6 — Fonctions trigonométriques
%% Tous les résultats ont été vérifiés par calcul formel (sympy).

\corrigesChapitre{6}{Fonctions trigonométriques}

\begin{deuxcolonnes}

\begin{corrige}{1}{Valeurs remarquables}
\begin{questions}
  \item $\frac12+\frac12=\res{1}$
  \item $1-1=\res{0}$
  \item Relation fondamentale : $\res{1}$
  \item $\dfrac{1/2}{\sqrt3/2}=\dfrac{1}{\sqrt3}=\res{\dfrac{\sqrt3}{3}}$, c'est $\tan\frac{\pi}{6}$.
  \item $\dfrac{\sqrt3}{2}\times\dfrac{\sqrt3}{2}=\res{\dfrac34}$
\end{questions}
\end{corrige}

\begin{corrige}{2}{Angles associés}
\begin{questions}
  \item $\cos\frac{5\pi}{6}=\cos\left(\pi-\frac{\pi}{6}\right)=\res{-\frac{\sqrt3}{2}}$\\[0.8mm]
    $\sin\frac{7\pi}{6}=\sin\left(\pi+\frac{\pi}{6}\right)=\res{-\frac12}$\\[0.8mm]
    $\cos\left(-\frac{2\pi}{3}\right)=\cos\left(\pi-\frac{\pi}{3}\right)=\res{-\frac12}$\\[0.8mm]
    $\sin\frac{3\pi}{4}=\sin\left(\pi-\frac{\pi}{4}\right)=\res{\frac{\sqrt2}{2}}$\\[0.8mm]
    $\cos\frac{11\pi}{6}=\cos\left(-\frac{\pi}{6}+2\pi\right)=\res{\frac{\sqrt3}{2}}$\\[0.8mm]
    $\sin\frac{17\pi}{4}=\sin\left(\frac{\pi}{4}+4\pi\right)=\res{\frac{\sqrt2}{2}}$
  \item $A(x)=-\cos x+\cos x=\res{0}$\\[0.8mm]
    $B(x)=\sin x+\sin x=\res{2\sin x}$\\[0.8mm]
    $C(x)=\cos x-\cos x+\cos x=\res{\cos x}$
  \item $\sin^{2}x=1-\frac{9}{25}=\frac{16}{25}$ et $\sin x\leq0$ sur $\intFF{-\frac{\pi}{2}}{0}$, donc $\res{\sin x=-\frac45}$.
\end{questions}
\end{corrige}

\begin{corrige}{3}{Du cercle à la courbe}
\begin{questions}
  \item Symétries par rapport aux axes : $\sin\frac{\pi}{6}=\sin\frac{5\pi}{6}=\res{\frac12}$ et
    $\sin\left(-\frac{\pi}{6}\right)=\res{-\frac12}$.
  \item $M(\pi-x)$ et $M(x)$ sont symétriques par rapport à $(Oy)$ : même ordonnée. Courbe symétrique par rapport à la
    \res{\text{droite } x=\tfrac{\pi}{2}}.
  \item L'ordonnée de $M(x)$ croît de $0$ à $1$, puis décroît : \res{\text{croissante sur } \intFF{0}{\frac{\pi}{2}}},
    \res{\text{décroissante sur } \intFF{\frac{\pi}{2}}{\pi}}.
  \item $\res{S=\ens{-\frac{\pi}{3}\pv\frac{\pi}{3}}}$ ; arc de gauche :
    $\res{S=\intFF{-\pi}{-\frac{\pi}{3}}\cup\intFF{\frac{\pi}{3}}{\pi}}$ (courbe sous $y=\frac12$).
\end{questions}
\end{corrige}

\begin{corrige}{4}{Parité et périodicité en images}
\begin{questions}
  \item $f(-x)=f(x)$, $g(-x)=(-x)\sin(-x)=g(x)$ : \res{f,\,g\text{ paires}} ; $k(-x)=-k(x)$ : \res{k\text{ impaire}} ;
    $h\left(\frac{\pi}{4}\right)=\sqrt2$, $h\left(-\frac{\pi}{4}\right)=0$ : \res{h\text{ ni paire ni impaire}}.
  \item $f(x+\pi)=\cos(2x+2\pi)=f(x)$ ; $k(x+\pi)=(-\sin x)(-\cos x)=k(x)$ ; $h(x+2\pi)=h(x)$.
  \item \otRond[otGreenDark]{1} $=\res{k}$ (centre $O$, maximum $\frac12$) ; \otRond[otGreenDark]{2} $=\res{h}$ (sans
    symétrie) ; \otRond[otGreenDark]{3} $=\res{g}$ (axe $(Oy)$, non bornée) ; \otRond[otGreenDark]{4} $=\res{f}$ (axe $(Oy)$,
    $f(0)=1$).
  \item $k$ est $\pi$-périodique et impaire : tracé sur $\intFF{0}{\frac{\pi}{2}}$, puis \res{\text{symétrie de centre } O}
    et \res{\text{translations de vecteurs } p\pi\,\vec\imath}.
\end{questions}
\end{corrige}

\begin{corrige}{5}{Dérivées}
\methode{$\big(\sin u\big)'=u'\cos u$ et $\big(\cos u\big)'=-u'\sin u$}
\begin{questions}
  \item $\res{f'(x)=3\cos(3x)}$
  \item $\res{g'(x)=2x\cos x-x^{2}\sin x}$
  \item $\res{h'(x)=\dfrac{2x}{\cos^{2}\left(x^{2}+1\right)}}$, là où $\cos\left(x^{2}+1\right)\neq0$.
  \item $k'(x)=\cos x\cos x+\sin x\,(-\sin x)$\\ $=\cos^{2}x-\sin^{2}x=\res{\cos(2x)}$
  \item $\res{m'(x)=\dfrac{x\cos x-\sin x}{x^{2}}}$
\end{questions}
\end{corrige}

\begin{corrige}{6}{Équations}
\begin{questions}
  \item $\sin x=\sin\frac{\pi}{6}$ : $x=\frac{\pi}{6}+2k\pi$ ou $x=\frac{5\pi}{6}+2k\pi$ ; $\res{S=\ens{\frac{\pi}{6}\pv\frac{5\pi}{6}}}$.
  \item $\cos x=\cos\frac{3\pi}{4}$ : $x=\pm\frac{3\pi}{4}+2k\pi$ ; $\res{S=\ens{\frac{3\pi}{4}\pv\frac{5\pi}{4}}}$.
  \item $x=\frac{\pi}{4}+k\pi$ : $\res{S=\ens{\frac{\pi}{4}\pv\frac{5\pi}{4}}}$.
  \item Avec $X=\sin x$ : $2X^{2}-3X+1=(X-1)(2X-1)$, donc $\sin x=1$ ou $\sin x=\frac12$ : $\res{S=\ens{\frac{\pi}{6}\pv\frac{\pi}{2}\pv\frac{5\pi}{6}}}$.
  \item L'équation s'écrit $2\sin x\cos x-\cos x=0$, soit $\cos x\,(2\sin x-1)=0$ : $\cos x=0$ ou $\sin x=\frac12$.\\
    $\res{S=\ens{\frac{\pi}{6}\pv\frac{\pi}{2}\pv\frac{5\pi}{6}\pv\frac{3\pi}{2}}}$
\end{questions}
\end{corrige}

\begin{corrige}{7}{Inéquations}
\begin{questions}
  \item $\sin x=\frac{\sqrt3}{2}$ en $\frac{\pi}{3}$ et $\frac{2\pi}{3}$ ; l'arc supérieur convient : $\res{S=\intFF{\frac{\pi}{3}}{\frac{2\pi}{3}}}$.
  \item $X=2x$ décrit $\intFF{0}{2\pi}$, et $\cos X<\frac12\iff X\in\intOO{\frac{\pi}{3}}{\frac{5\pi}{3}}$ :
    $\res{S=\intOO{\frac{\pi}{6}}{\frac{5\pi}{6}}}$.
  \item $\tan$ est croissante sur chaque intervalle de définition et vaut $1$ en $\frac{\pi}{4}$ et $\frac{5\pi}{4}$ :
    $\res{S=\intFF{0}{\frac{\pi}{4}}\cup\intOF{\frac{\pi}{2}}{\frac{5\pi}{4}}\cup\intOO{\frac{3\pi}{2}}{2\pi}}$.
  \item $\sqrt2\sin\left(x+\frac{\pi}{4}\right)$\\
    $=\sqrt2\left(\sin x\cos\frac{\pi}{4}+\cos x\sin\frac{\pi}{4}\right)$\\
    $=\sqrt2\times\frac{\sqrt2}{2}(\sin x+\cos x)=\sin x+\cos x$.\\
    Le maximum est $\res{\sqrt2}$, atteint lorsque $x+\frac{\pi}{4}=\frac{\pi}{2}$, soit en $x=\frac{\pi}{4}$.\\
    Enfin, l'inéquation équivaut à $\sin X\geq\frac{\sqrt2}{2}$ avec $X=x+\frac{\pi}{4}\in\intFF{\frac{\pi}{4}}{\frac{9\pi}{4}}$, soit
    $X\in\intFF{\frac{\pi}{4}}{\frac{3\pi}{4}}$ ou $X=\frac{9\pi}{4}$ : $\res{S=\intFF{0}{\frac{\pi}{2}}\cup\ens{2\pi}}$.
\end{questions}
\end{corrige}

\begin{corrige}{8}{Formules et identités}
\begin{questions}
  \item Les termes $\cos a\sin b$ s'éliminent : $\res{2\sin a\cos b}$.
  \item $\cos\frac{\pi}{12}=\cos\frac{\pi}{3}\cos\frac{\pi}{4}+\sin\frac{\pi}{3}\sin\frac{\pi}{4}$\\
    $=\frac12\times\frac{\sqrt2}{2}+\frac{\sqrt3}{2}\times\frac{\sqrt2}{2}=\res{\dfrac{\sqrt2+\sqrt6}{4}}$
  \item $\cos(2x)=1-2\sin^{2}x$ donne $\res{\sin^{2}x=\dfrac{1-\cos(2x)}{2}}$.
  \item $\dfrac{2\sin x\cos x}{2\cos^{2}x}=\dfrac{\sin x}{\cos x}=\res{\tan x}$
  \item $\sin\frac{5\pi}{12}=\sin\left(\frac{\pi}{4}+\frac{\pi}{6}\right)$\\
    $=\frac{\sqrt2}{2}\times\frac{\sqrt3}{2}+\frac{\sqrt2}{2}\times\frac12=\res{\dfrac{\sqrt6+\sqrt2}{4}}$
\end{questions}
\end{corrige}

\begin{corrige}{9}{Limites}
\begin{questions}
  \item $\res{1}$ (limite de référence).
  \item $\dfrac{1-\cos x}{x^{2}}=\dfrac{1-\cos^{2}x}{x^{2}(1+\cos x)}$\\[1mm]
    $=\left(\dfrac{\sin x}{x}\right)^{2}\times\dfrac{1}{1+\cos x}\to1\times\frac12=\res{\frac12}$
  \item À gauche : $\tan x\to+\infty$ et $x-\frac{\pi}{2}\to0^{-}$ ; à droite : $\tan x\to-\infty$ et $x-\frac{\pi}{2}\to0^{+}$.
    Dans les deux cas, la limite est $\res{-\infty}$.
  \item Oui : $\lim_{0}f=1$, donc on prolonge $f$ en posant $\res{f(0)=1}$.
  \item Pour $x>0$, $-\frac1x\leq\frac{\sin x}{x}\leq\frac1x$ : par encadrement, la limite est $\res{0}$.
\end{questions}
\end{corrige}

\end{deuxcolonnes}

%% ---------------------------------------------------------------------
\begin{corrige}{10}{Études de fonctions}
\noindent
\begin{minipage}[t]{0.52\linewidth}\vspace{0pt}\raggedright
\begin{questions}
  \item $f'(x)=\cos x-\sin x$. Or $\cos x=\sin x\iff\cos x=\cos\left(\frac{\pi}{2}-x\right)\iff x=\frac{\pi}{4}+k\pi$ : sur
    $\intFF{0}{2\pi}$, $f'$ s'annule en $\frac{\pi}{4}$ et $\frac{5\pi}{4}$. Comme $f'(0)=1$, $f'(\pi)=-1$ et $f'(2\pi)=1$, on
    obtient le tableau ci-contre : $\res{\max f=\sqrt2}$ et $\res{\min f=-\sqrt2}$.
\end{questions}
\end{minipage}\hfill
\begin{minipage}[t]{0.44\linewidth}\vspace{1mm}
\centering
\begin{tableau}
  \otTab[1.15]{$x$/0.8,$f'(x)$/0.8,$f$/1.6}{$0$,$\frac{\pi}{4}$,$\frac{5\pi}{4}$,$2\pi$}
  \tkzTabLine{,+,z,-,z,+,}
  \tkzTabVar{-/$1$,+/$\sqrt2$,-/$-\sqrt2$,+/$1$}
\end{tableau}
\end{minipage}

\noindent
\begin{minipage}[t]{0.42\linewidth}\vspace{0pt}\raggedright
\begin{questions}[start=2]
  \item $g(x)=-\left(\cos^{2}x-\sin^{2}x\right)=-\cos(2x)$, donc $g'(x)=2\sin(2x)$, qui s'annule en $0$, $\frac{\pi}{2}$, $\pi$,
    $\frac{3\pi}{2}$, $2\pi$. $\res{\max g=1}$ en $\frac{\pi}{2}$ et $\frac{3\pi}{2}$ ; $\res{\min g=-1}$ en $0$, $\pi$, $2\pi$.
\end{questions}
\end{minipage}\hfill
\begin{minipage}[t]{0.54\linewidth}\vspace{1mm}
\centering
\begin{tableau}
  \otTab[1.1]{$x$/0.8,$g'(x)$/0.8,$g$/1.6}{$0$,$\frac{\pi}{2}$,$\pi$,$\frac{3\pi}{2}$,$2\pi$}
  \tkzTabLine{z,+,z,-,z,+,z,-,z}
  \tkzTabVar{-/$-1$,+/$1$,-/$-1$,+/$1$,-/$-1$}
\end{tableau}
\end{minipage}

\noindent
\begin{minipage}[t]{0.52\linewidth}\vspace{0pt}\raggedright
\begin{questions}[start=3]
  \item $h'(x)=\dfrac{\cos x\,(1+\cos x)+\sin^{2}x}{(1+\cos x)^{2}}=\dfrac{1+\cos x}{(1+\cos x)^{2}}=\res{\dfrac{1}{1+\cos x}}>0$ :
    $h$ est croissante sur $\intFO{0}{\pi}$ et sur $\intOO{\pi}{2\pi}$. En $\pi$, $1+\cos x\to0^{+}$ et $\sin x\to0^{\pm}$ :
    $\res{\lim_{\pi^{-}}h=+\infty}$ et $\res{\lim_{\pi^{+}}h=-\infty}$. La droite $x=\pi$ est asymptote.
\end{questions}
\end{minipage}\hfill
\begin{minipage}[t]{0.44\linewidth}\vspace{1mm}
\centering
\begin{tableau}
  \otTab[1.6]{$x$/0.8,$h'(x)$/0.8,$h$/1.6}{$0$,$\pi$,$2\pi$}
  \tkzTabLine{,+,d,+,}
  \tkzTabVar{-/$0$,+D-/$+\infty$/$-\infty$,+/$0$}
\end{tableau}
\end{minipage}

\begin{questions}[start=4]
  \item $k(x)=\frac12\sin(2x)$, donc $k'(x)=\cos(2x)$, qui s'annule en $\frac{\pi}{4}$, $\frac{3\pi}{4}$, $\frac{5\pi}{4}$, $\frac{7\pi}{4}$.
    $k$ croît sur $\intFF{0}{\frac{\pi}{4}}$, décroît sur $\intFF{\frac{\pi}{4}}{\frac{3\pi}{4}}$, puis recommence ($k$ est $\pi$-périodique).
    $\res{\max k=\frac12}$ en $\frac{\pi}{4}$ et $\frac{5\pi}{4}$ ; $\res{\min k=-\frac12}$ en $\frac{3\pi}{4}$ et $\frac{7\pi}{4}$.
  \item $m'(x)=\dfrac{1}{\cos^{2}x}-1=\dfrac{1-\cos^{2}x}{\cos^{2}x}=\res{\tan^{2}x}\geq0$ : $m$ est croissante, avec $m(0)=0$ et
    $\lim_{\frac{\pi}{2}^{-}}m=+\infty$. Donc $m(x)\geq m(0)=0$ : $\res{\tan x\geq x}$ sur $\intFO{0}{\frac{\pi}{2}}$.
\end{questions}
\end{corrige}

%% ---------------------------------------------------------------------
\begin{corrige}{11}{Synthèse}
\begin{questions}
  \item $\sin(3x)=\sin x\iff3x=x+2k\pi$ ou $3x=\pi-x+2k\pi\iff x=k\pi$ ou $x=\frac{\pi}{4}+k\frac{\pi}{2}$ :
    $\res{S=\ens{0\pv\frac{\pi}{4}\pv\frac{3\pi}{4}\pv\pi\pv\frac{5\pi}{4}\pv\frac{7\pi}{4}\pv2\pi}}$.
\end{questions}
\noindent
\begin{minipage}[t]{0.5\linewidth}\vspace{0pt}\raggedright
\begin{questions}[start=2]
  \item $f'(x)=\dfrac{\cos x\,(1+\sin x)-\sin x\cos x}{(1+\sin x)^{2}}=\res{\dfrac{\cos x}{(1+\sin x)^{2}}}$, du signe de $\cos x$.
    $\res{\max f=f\left(\frac{\pi}{2}\right)=\frac12}$. En $\frac{3\pi}{2}$, $\sin x\to-1$ et $1+\sin x\to0^{+}$ : la limite est
    $-\infty$ des deux côtés ; la droite $x=\frac{3\pi}{2}$ est asymptote.
\end{questions}
\end{minipage}\hfill
\begin{minipage}[t]{0.46\linewidth}\vspace{1mm}
\centering
\begin{tableau}
  \otTab[1.25]{$x$/0.8,$f'(x)$/0.8,$f$/1.6}{$0$,$\frac{\pi}{2}$,$\frac{3\pi}{2}$,$2\pi$}
  \tkzTabLine{,+,z,-,d,+,}
  \tkzTabVar{-/$0$,+/$\frac12$,-D-/$-\infty$/$-\infty$,+/$0$}
\end{tableau}
\end{minipage}
\begin{questions}[start=3]
  \item $X=2x$ décrit $\intFF{0}{4\pi}$ et $\cos X\geq\frac12\iff X\in\intFF{0}{\frac{\pi}{3}}\cup\intFF{\frac{5\pi}{3}}{\frac{7\pi}{3}}\cup\intFF{\frac{11\pi}{3}}{4\pi}$.\\
    Donc $\res{S=\intFF{0}{\frac{\pi}{6}}\cup\intFF{\frac{5\pi}{6}}{\frac{7\pi}{6}}\cup\intFF{\frac{11\pi}{6}}{2\pi}}$.
  \item $g'(x)=2x\cos\left(x^{2}\right)$, puis $\res{g''(x)=2\cos\left(x^{2}\right)-4x^{2}\sin\left(x^{2}\right)}$.
  \item $2x=\frac{\pi}{3}+k\pi$, soit $x=\frac{\pi}{6}+k\frac{\pi}{2}$ : $\res{S=\ens{\frac{\pi}{6}\pv\frac{2\pi}{3}}}$.
\end{questions}
\end{corrige}

%% ---------------------------------------------------------------------
\begin{corrige}{12}{Suites et trigonométrie}
\begin{questions}
  \item $u_0=\cos\pi=\res{-1}$, $u_1=\cos\frac{\pi}{2}=\res{0}$, $u_2=\cos\frac{\pi}{4}=\res{\frac{\sqrt2}{2}}$. Comme $\frac{\pi}{2^{n}}\to0$ et que
    $\cos$ est continue en $0$, $\res{\lim u_n=\cos0=1}$.
  \item $v_{n+8}=\sin\left(\frac{n\pi}{4}+2\pi\right)=v_n$. Les valeurs $v_0,\dots,v_7$ sont $0$, $\frac{\sqrt2}{2}$, $1$, $\frac{\sqrt2}{2}$, $0$,
    $-\frac{\sqrt2}{2}$, $-1$, $-\frac{\sqrt2}{2}$ : aucune période plus courte ne convient, la suite est $\res{\text{périodique de période }8}$.
  \item $w_2=\frac12\left(\sin\frac{\pi}{2}+\sin\pi\right)=\res{\frac12}$ et $w_4=\frac14\left(\frac{\sqrt2}{2}+1+\frac{\sqrt2}{2}+0\right)=\res{\frac{1+\sqrt2}{4}}$.
  \item $u_1=\res{1}$, donc $u_{n+1}=u_n-u_{n-1}$ et $u_2=\res{-1}$. On démontre par récurrence la propriété
    $\mathcal{P}(n)$ : « $u_n=2\cos\frac{n\pi}{3}$ et $u_{n+1}=2\cos\frac{(n+1)\pi}{3}$ ».\\
    \etape{Initialisation} $u_0=2=2\cos0$ et $u_1=1=2\cos\frac{\pi}{3}$.\\
    \etape{Hérédité} Si $\mathcal{P}(n)$ est vraie, $u_{n+2}=2\cos\frac{\pi}{3}\times2\cos\frac{(n+1)\pi}{3}-2\cos\frac{n\pi}{3}$, d'où,
    avec la formule de linéarisation :
    \[ u_{n+2}=2\left[\cos\frac{(n+2)\pi}{3}+\cos\frac{n\pi}{3}\right]-2\cos\frac{n\pi}{3}=2\cos\frac{(n+2)\pi}{3}, \]
    donc $\mathcal{P}(n+1)$ est vraie.\\
    \etape{Conclusion} $\res{u_n=2\cos\frac{n\pi}{3}}$ pour tout $n\in\mathbb{N}$.
  \item Les valeurs de $\cos\frac{k\pi}{3}$ pour $k=0,\dots,6$ sont $1$, $\frac12$, $-\frac12$, $-1$, $-\frac12$, $\frac12$, $1$ :
    $t_1=\res{\frac32}$, $t_2=\res{1}$, $t_6=\res{1}$. La somme de six termes consécutifs est nulle, car
    $\cos\frac{(k+3)\pi}{3}=-\cos\frac{k\pi}{3}$ ; donc $t_{n+6}=t_n$ : $\res{(t_n)\text{ est périodique de période }6}$.
\end{questions}
\end{corrige}

%% ---------------------------------------------------------------------
\begin{corrige}{13}{Étude complète d'une fonction trigonométrique}
\partieExo{A}{Réduction de l'intervalle d'étude}
\begin{questions}
  \item $\sin$ est $2\pi$-périodique et $\sin(2x+4\pi)=\sin(2x)$, donc $f(x+2\pi)=f(x)$. De plus
    $f(-x)=-2\sin x-\sin(2x)=-f(x)$ : $f$ est impaire. $\Cf$ est invariante par la $\res{\text{translation de vecteur }2\pi\,\vec\imath}$
    et par la $\res{\text{symétrie de centre }O}$.
  \item Par périodicité, on étudie $f$ sur $\intFF{-\pi}{\pi}$ ; par imparité, sur $\intFF{0}{\pi}$, puis on complète par symétrie.
\end{questions}
\partieExo{B}{Variations}
\noindent
\begin{minipage}[t]{0.54\linewidth}\vspace{0pt}\raggedright
\begin{questions}[resume]
  \item $f'(x)=2\cos x+2\cos(2x)=4\cos^{2}x+2\cos x-2=2\left(2\cos^{2}x+\cos x-1\right)$, et
    $(2\cos x-1)(\cos x+1)=2\cos^{2}x+\cos x-1$.
  \item Sur $\intFF{0}{\pi}$ : $\cos x+1\geq0$ (nul seulement en $\pi$) et $2\cos x-1>0\iff\cos x>\frac12\iff x\in\intFO{0}{\frac{\pi}{3}}$.
    $\res{\max f=f\left(\frac{\pi}{3}\right)=\frac{3\sqrt3}{2}\approx2{,}60}$.
\end{questions}
\end{minipage}\hfill
\begin{minipage}[t]{0.42\linewidth}\vspace{1mm}
\centering
\begin{tableau}
  \otTab[1.6]{$x$/0.8,$f'(x)$/0.8,$f$/1.6}{$0$,$\frac{\pi}{3}$,$\pi$}
  \tkzTabLine{,+,z,-,z}
  \tkzTabVar{-/$0$,+/$\frac{3\sqrt3}{2}$,-/$0$}
\end{tableau}
\end{minipage}
\partieExo{C}{Une équation et une tangente}
\noindent
\begin{minipage}[t]{0.5\linewidth}\vspace{0pt}\raggedright
\begin{questions}[start=5]
  \item $f\left(\frac{\pi}{2}\right)=2+\sin\pi=2$. Sur $\intFF{0}{\frac{\pi}{3}}$, $f$ est continue, strictement croissante, et
    $f(0)=0<2<\frac{3\sqrt3}{2}$ : une unique solution $\alpha$. Sur $\intFF{\frac{\pi}{3}}{\pi}$, $f$ est strictement décroissante :
    une unique solution, qui est $\frac{\pi}{2}$. $\res{\text{Deux solutions}}$ ; comme $f(0{,}57)\approx1{,}988$ et
    $f(0{,}58)\approx2{,}013$, $\res{0{,}57<\alpha<0{,}58}$.
  \item $f(0)=0$ et $f'(0)=4$ : $\res{T:y=4x}$. Puis $f'(x)-4=4\cos^{2}x+2\cos x-6=2\left(2\cos^{2}x+\cos x-3\right)$
    $=2(\cos x-1)(2\cos x+3)$. Or $\cos x-1\leq0$ et $2\cos x+3>0$ : la fonction $\varphi:x\mapsto f(x)-4x$ est
    décroissante sur $\R$, et $\varphi(0)=0$. Donc $\varphi(x)\leq0$ pour $x\geq0$ : $\res{\Cf\text{ est sous }T}$ sur $\intFO{0}{+\infty}$.
  \item Voir la figure ci-contre.
\end{questions}
\end{minipage}\hfill
\begin{minipage}[t]{0.46\linewidth}\vspace{2mm}
\centering
\begin{tikzpicture}
\begin{axis}[otGraphe, width=72mm, xmin=-3.5, xmax=3.7, ymin=-3.1, ymax=3.3,
  xtick={-3.1416,-1.5708,1.0472,1.5708,3.1416}, xticklabels={$-\pi$,$-\frac{\pi}{2}$,$\frac{\pi}{3}$,$\frac{\pi}{2}$,$\pi$},
  ytick={-2,-1,1,2}]
  \draw[otMuted,line width=0.5pt,dash pattern=on 2pt off 1.5pt] (axis cs:-0.2,2) -- (axis cs:2.4,2);
  \addplot[otGreen,line width=0.9pt,domain=-0.77:0.79] {4*x};
  \addplot[otNavy,line width=1.2pt,domain=-3.1416:3.1416,samples=200] {2*sin(deg(x))+sin(deg(2*x))};
  \addplot[only marks,mark=*,mark size=1.4pt,otNavy] coordinates {(1.0472,2.598) (0.5748,2) (1.5708,2)};
  \node[anchor=south,otEtiquette,font=\scriptsize] at (axis cs:1.05,2.62) {$\frac{3\sqrt3}{2}$};
  \node[anchor=south east,font=\scriptsize,inner sep=1.5pt] at (axis cs:0.55,2.03) {$\alpha$};
  \node[anchor=west,font=\scriptsize,text=otGreenDark] at (axis cs:0.8,3.1) {$T$};
  \node[anchor=north east,font=\scriptsize,text=otNavy] at (axis cs:-2.4,-0.9) {$\Cf$};
\end{axis}
\end{tikzpicture}\\[0.5mm]
\legende{$\Cf$ est symétrique par rapport à l'origine.}
\end{minipage}
\end{corrige}

%% ---------------------------------------------------------------------
\begin{corrige}{14}{Oscillations amorties}
\partieExo{A}{Enveloppes et limite}
\begin{questions}
  \item $-1\leq\cos(2t)\leq1$ et $e^{-t/2}>0$ : on multiplie par $e^{-t/2}$. Comme $e^{-t/2}$ tend vers $0$ lorsque $t$ tend
    vers $+\infty$, le théorème des gendarmes donne $\res{\lim_{t\to+\infty}x(t)=0}$ : le ressort revient vers sa position d'équilibre.
  \item $e^{-t/2}\neq0$, donc $x(t)=0\iff\cos(2t)=0\iff2t=\frac{\pi}{2}+k\pi$ : $\res{t=\frac{\pi}{4}+k\frac{\pi}{2}}$, $k\in\mathbb{N}$.
  \item $x(t)=e^{-t/2}\iff\cos(2t)=1\iff2t=2k\pi$ : $\res{t=k\pi}$, $k\in\mathbb{N}$.
\end{questions}
\partieExo{B}{Vitesse et amortissement}
\noindent
\begin{minipage}[t]{0.5\linewidth}\vspace{0pt}\raggedright
\begin{questions}[resume]
  \item $x'(t)=-\frac12e^{-t/2}\cos(2t)+e^{-t/2}\times\left(-2\sin(2t)\right)$, d'où la forme demandée. En $t=k\pi$ :
    $\cos(2k\pi)=1$ et $\sin(2k\pi)=0$, donc $x'(k\pi)=-\frac12e^{-k\pi/2}$, qui est aussi le nombre dérivé de $t\mapsto e^{-t/2}$
    en $k\pi$. Mêmes valeurs, mêmes nombres dérivés : $\res{\text{même tangente}}$, la courbe « touche » son enveloppe.
  \item $x(t+\pi)=e^{-t/2}e^{-\pi/2}\cos(2t+2\pi)=\res{e^{-\pi/2}\,x(t)}$, avec $e^{-\pi/2}\approx0{,}21$ : toutes les $\pi\approx3{,}14$
    secondes, l'amplitude est multipliée par $0{,}21$ environ.
  \item $e^{-t/2}<0{,}01\iff-\frac t2<\ln0{,}01\iff t>2\ln100=4\ln10$ : à partir de $\res{t\approx9{,}21\text{ s}}$.
\end{questions}
\end{minipage}\hfill
\begin{minipage}[t]{0.46\linewidth}\vspace{1mm}
\centering
\begin{tikzpicture}
\begin{axis}[otSchema, axis lines=middle, width=74mm, height=50mm, xmin=0, xmax=10.6, ymin=-1.15, ymax=1.2,
  xtick={3.1416,6.2832,9.4248}, xticklabels={$\pi$,$2\pi$,$3\pi$}, ytick={-1,1},
  xlabel={$t$}, xlabel style={at={(1,0.5)},anchor=west},
  grid=major, grid style={otRule!70,line width=0.3pt}]
  \addplot[otGreen,dashed,line width=0.7pt,domain=0:10.4,samples=60] {exp(-x/2)};
  \addplot[otGreen,dashed,line width=0.7pt,domain=0:10.4,samples=60] {-exp(-x/2)};
  \addplot[otNavy,line width=1.1pt,domain=0:10.4,samples=300] {exp(-x/2)*cos(deg(2*x))};
  \addplot[only marks,mark=*,mark size=1.3pt,otNavy] coordinates {(0,1) (3.1416,0.2079) (6.2832,0.0432)};
  \node[anchor=west,font=\scriptsize,text=otGreenDark] at (axis cs:1.15,0.62) {$e^{-t/2}$};
\end{axis}
\end{tikzpicture}\\[0.5mm]
\legende{La courbe oscille entre ses deux enveloppes et les touche en $t=k\pi$.}
\end{minipage}
\end{corrige}

%% ---------------------------------------------------------------------
\begin{corrige}{15}{Encadrements de sinus et cosinus}
\noindent
\begin{minipage}[t]{0.54\linewidth}\vspace{0pt}\raggedright
\partieExo{A}{Trois inégalités}
\begin{questions}
  \item $g'(x)=1-\cos x\geq0$ : $g$ est croissante sur $\intFO{0}{+\infty}$ et $g(0)=0$, donc $g(x)\geq0$ : $\res{\sin x\leq x}$.
  \item $h'(x)=-\sin x+x=g(x)\geq0$ : $h$ est croissante et $h(0)=0$, donc $h(x)\geq0$ : $\res{\cos x\geq1-\frac{x^{2}}{2}}$.
  \item $k(x)=\sin x-x+\frac{x^{3}}{6}$ a pour dérivée $k'(x)=\cos x-1+\frac{x^{2}}{2}=h(x)\geq0$, et $k(0)=0$ : $\res{\sin x\geq x-\frac{x^{3}}{6}}$.
\end{questions}
\partieExo{B}{Applications}
\begin{questions}[resume]
  \item $0{,}1-\frac{0{,}001}{6}\leq\sin(0{,}1)\leq0{,}1$, soit $\res{0{,}099\,83\leq\sin(0{,}1)\leq0{,}1}$, d'amplitude
    $\frac{1}{6\,000}\approx1{,}7\times10^{-4}$ (en réalité, $\sin(0{,}1)\approx0{,}099\,833$). De même, $\res{\cos(0{,}2)\geq0{,}98}$.
  \item On divise l'encadrement $x-\frac{x^{3}}{6}\leq\sin x\leq x$ par $x>0$. Les deux bornes tendent vers $1$ en $0$ : par
    encadrement, $\lim_{0^{+}}\frac{\sin x}{x}=1$. La fonction $x\mapsto\frac{\sin x}{x}$ étant paire, la limite à gauche vaut aussi
    $1$ : $\res{\lim_{x\to0}\frac{\sin x}{x}=1}$.
\end{questions}
\end{minipage}\hfill
\begin{minipage}[t]{0.42\linewidth}\vspace{2mm}
\centering
\begin{tikzpicture}
\begin{axis}[otGraphe, width=66mm, xmin=-0.2, xmax=3.5, ymin=-0.6, ymax=3.1,
  xtick={1,2,3}, ytick={1,2,3}]
  \addplot[otGreen,dashed,line width=0.8pt,domain=0:3] {x};
  \addplot[otGreenDark,line width=0.8pt,domain=0:2.66,samples=80] {x-x^3/6};
  \addplot[otNavy,line width=1.2pt,domain=0:3.3,samples=120] {sin(deg(x))};
  \node[anchor=north west,font=\scriptsize,text=otGreenDark] at (axis cs:2.4,2.1) {$y=x$};
  \node[anchor=west,font=\scriptsize,text=otNavy] at (axis cs:2.55,0.75) {$y=\sin x$};
  \node[anchor=east,font=\scriptsize,text=otGreenDark] at (axis cs:2.0,0.3) {$y=x-\frac{x^{3}}{6}$};
\end{axis}
\end{tikzpicture}\\[0.5mm]
\legende{Près de $0$, $\sin x$ est très bien approché par $x-\frac{x^{3}}{6}$.}
\end{minipage}
\end{corrige}
